\section{Lower degree from the ordinary strong equation}
\label{sec:ordinary-new}

The inputless arithmetic of packages~06 and~08 can be combined with the
ordinary strong equation used in package~02
\cite{src02,src06,src08}. This improves degree while preserving the
fifteen positive unknowns and the four fixed coefficient ports. Keep
the previously defined polynomials
\(q,C,X,Y,E,k,c,a,H,\Delta,D,R,V\), and keep the first, main, index and
transport factors \(N_0,N_1,N_k,N_t\). Replace the strong and auxiliary
block by the polynomial abbreviations
\begin{equation}\label{eq:ordinary-data}
 I=ic^2,\qquad W=\Delta(f^2-1),\qquad
 N_a=W(V^2-y^2)+y^2.
\end{equation}
Here \(I,W\) are computed expressions, not additional unknowns. In
particular, the strong equation is \(I^2=W\), rather than the
normalized equation involving \((i\Delta c^2)^2\).

\begin{center}
\begin{tabular}{@{}lrrrr@{}}
\toprule
Construction & Unknowns & Equations & Operations & Degree\\
\midrule
Ordinary system & 15 & 6 & \(34M+27A=61\) & 28\\
Grouped ordinary system & 15 & 4 & \(36M+27A=63\) & 28\\
Four-square polynomial & 15 & 1 & \(40M+34A=74\) & 56\\
Ordinary product & 15 & 1 & \(39M+30A=69\) & 99\\
Quadratic lift & 33 & 24 & not optimized & 2\\
Its sum of squares & 33 & 1 & not optimized & 4\\
\bottomrule
\end{tabular}
\end{center}
All unknown counts in this table use the strictly positive domain.
The operation model is unchanged: each binary addition, subtraction
or multiplication costs one, including multiplication by a numeral.

\subsection{Six equations, four equations, and one polynomial}

The ordinary six-equation system is
\begin{equation}\label{eq:ordinary-six}
 \boxed{N_0=N_1=N_a=N_k=N_t=1,\qquad I^2=W.}
\end{equation}
Its literal circuit uses 61 operations, and its residual degrees, in
the displayed order, are
\[
 22,\quad18,\quad28,\quad7,\quad2,\quad22.
\]
The auxiliary coefficient \(W=\Delta(f^2-1)\) is intentionally used
in~\eqref{eq:ordinary-data}. Replacing it off the zero set by its
equal-on-solutions expression \(I^2\) would increase the degree of
the auxiliary residual.

\begin{proposition}[Grouping the unit factors]\label{prop:group-ordinary}
Over positive integers, system~\eqref{eq:ordinary-six} is tuplewise
equivalent to the four equations
\begin{equation}\label{eq:ordinary-four}
 \boxed{N_0N_t=1,\qquad N_1N_k=1,\qquad N_a=1,\qquad I^2=W.}
\end{equation}
\end{proposition}
\begin{proof}
Each factor in either product equation is an integer unit. The first
factor has the form
\(\tau^2-U(U+1)k^2\), where \(U=XY^2>1\); the
consecutive-product negative Pell equation excludes the value \(-1\).
Also \(\Delta=(a+2)^2-1\) is \(0\) or \(3\pmod4\), so
\(N_1=D^2-\Delta c^2\) cannot be \(-1\pmod4\). Thus
\(N_0=N_1=1\), and the product equations give \(N_t=N_k=1\).
The converse is immediate.
\end{proof}

The grouped residuals have exact degrees \(24,25,28,22\). Therefore
the particularly transparent single equation
\begin{equation}\label{eq:ordinary-sos74}
 \boxed{\begin{aligned}
 Q_{56}={}&(N_0N_t-1)^2+(N_1N_k-1)^2\\
         &+(N_a-1)^2+(I^2-W)^2=0
 \end{aligned}}
\end{equation}
uses 74 operations and has degree 56. This is four operations fewer
than the direct six-square finalizer, at the same degree and unknown
count. The degree is exact uniformly in the permitted fixed constants.
Write \(Q_0=\beta J\) and \(k_0=\eta+\zeta\) for highest homogeneous
forms. The unique degree-28 residual, \(N_a-1\), has highest form
\[
 w^2k_0^2s^4Q_0^{14}T^2f^4.
\]
Consequently the degree-56 part of \(Q_{56}\) is
\(w^4k_0^4s^8Q_0^{28}T^4f^8\), which is nonzero since \(\beta>0\).

\subsection{Why the ordinary strong equation is sufficient}

This reduction requires a different proof order from a purely normalized
strong-norm argument. One cannot simply divide the supplied \(i\) by
\(\Delta\). The relevant ordinary strong-rank theorem is the one proved
in the repository's relaxed auxiliary note~\cite{relaxed} and used in
package~02~\cite{src02}.

First consider a positive solution of~\eqref{eq:ordinary-six}. Its exact
transport equation forces \(C>0\). Indeed, if \(C\le0\), then
\[
 N_t=(K+w)C+1-F-(z-1)(q-1)\le0,
\]
contrary to \(N_t=1\). Thus \(F+C=q-\alpha<q\). The native mask bounds
and the same packing identity as in the inputless construction give
\begin{equation}\label{eq:ordinary-bounds}
 (2q-1)(q^2-1)<R<q^4-q^3,\qquad
 E\ge q^4>R+q^3,\qquad a>R+2.
\end{equation}
Set \(A=a+2\) and \(P_0=2XY^2+1\). The positive first and main norms
supply indices \(n,p\ge1\) such that
\[
 k=2\psi_{P_0}(n),\qquad c=\psi_A(p),\qquad D=\chi_A(p).
\]
Since \(P_0>A\) and \(c>k\), parameter monotonicity gives \(p>n\).
The first-index equation and \(P_0\equiv1\pmod E\) imply
\(2n\equiv R+1\pmod E\). Since \(0<R+1<E\), one has
\(n\ge(R+1)/2\). Elementary Pell growth therefore proves, before any
auxiliary argument,
\begin{equation}\label{eq:ordinary-large-c}
 c>A\Delta^2,\qquad c>2p,\qquad c>2R.
\end{equation}

Now \(I^2=W\) implies
\[
 f^2=1+\frac{i^2c^4}{\Delta},\qquad
 f>\frac{c^2}{A}>2c.
\]
Together with~\eqref{eq:ordinary-bounds}--\eqref{eq:ordinary-large-c},
this yields
\begin{equation}\label{eq:ordinary-gap}
 W-(Rf^2+c)=(\Delta-R)f^2-\Delta-c>0.
\end{equation}
The auxiliary equation is \(WV^2-(W-1)y^2=1\). If \(|V|>1\), then
\(V^2-1=(W-1)(y^2-V^2)\) gives \(|V|\ge2W-1\). But
\(V=c(Tf-1)-Rf^2>-(Rf^2+c)>-W\), excluding negative such values.
The values \(0,1,-1\) are excluded by the norm or by
\(V\equiv-c\pmod f\) and \(2<c<f/2\). Hence \(V>1\).

The relaxed strong-rank theorem now gives an index \(m\) with
\begin{equation}\label{eq:ordinary-rank}
 f=\chi_A(m),\quad p\mid m,\quad c\mid m,\quad
 ic^2=\Delta\psi_A(m),\quad m\ge c>2p.
\end{equation}
In particular \(c^2\mid f^2-1\), since
\(f^2-1=ic^2\psi_A(m)\). The proof-only integers
\[
 o=cT-Rf,\qquad j=Tf-R\frac{f^2-1}{c}-1
\]
are integral and positive: \(of=V+c>0\) and \(jc=V+R>0\).
They restore precisely \(V=of-c=jc-R\).

The retained odd auxiliary-index and signed step-down argument therefore
applies using~\eqref{eq:ordinary-rank}; it needs \(c\mid m\), not the
stronger and unjustified assertion \(c^2\mid\psi_A(m)\). It gives
\(R\equiv\pm p\pmod c\). Since \(0<R,p<c/2\), this forces \(p=R\).
The no-wrap argument gives \(2n=R+1\). The lower ratio estimate and main
projection then recover \(X=2^R\), followed by \(q=2^t\), the retained
\(q^3\)-divisibility of \(Y\), and the half-binomial population
threshold. The native masks, synchronization and forced Start now decode
the finite accepting computation exactly as in the inputless proof
\cite{src06,src08}.

For completeness, take the canonical normalized solution of the inputless
construction and replace only
\[
 i_{\mathrm{ordinary}}=\Delta i_{\mathrm{normalized}}.
\]
Its strong equality becomes \(I^2=W\), and its auxiliary coefficient is
unchanged. This is an existence-preserving map in the required direction;
no division by \(\Delta\) or coordinatewise equivalence of arbitrary
strong-norm tuples is asserted.

\subsection{A 69-operation product of degree 99}

The ordinary unit factor is \(N_s^{\mathrm{ord}}=1+I^2-W\). The equation
\begin{equation}\label{eq:ordinary-product69}
 \boxed{P_{69}=N_0N_1N_aN_kN_tN_s^{\mathrm{ord}}-1=0}
\end{equation}
has the same positive zero set as~\eqref{eq:ordinary-six}. It costs
69 operations and has degree 99. Its factor degrees are
\(22,18,28,7,2,22\), and its nonzero leading form is
\[
 8\gamma h i^2w^7s^{16}k_0^9Q_0^{55}T^2f^4
       \bigl(w(Q_0-F-\alpha)-zQ_0\bigr).
\]

Here the signed bootstrap must be stated separately. A zero makes every
factor a unit. The first and main norm signs are positive as before.
The factor \(N_s^{\mathrm{ord}}\) cannot be \(-1\pmod4\): when \(\Delta\equiv0\),
it is \(1+I^2\), and when \(\Delta\equiv3\), it is \(I^2+f^2\).
Thus \(I^2=W\), and the square coefficient in \(N_a\) also excludes
its negative sign. Write \(N_k=\varepsilon\), \(N_t=\nu\), where
\(\varepsilon,\nu\in\{1,-1\}\) and \(\varepsilon\nu=1\).

Transport first gives only \(C\ge0\). Its exceptional zero case is
\(C=0,F=2,z=1,\nu=-1\); retain it. The lower packing bound is still
\((2q-1)(q^2-1)<R\): for \(C>0\), use \(F+C\le q-1\), while in
the exceptional case \(q^2-C-qF=q^2-2q>2q-1\). For the upper bound,
the weaker estimate \(q^2-C-qF\le q^2-q\) suffices and gives
\[
 R<q^4-q^3+q^2-1,\qquad
 E\ge q^4>R+q^3-q^2+1>R+5.
\]
Now \(n\ge(R-1)/2\), \(p>n\), and \(k=R+\varepsilon+hE>2R\)
give~\eqref{eq:ordinary-large-c}. The ordinary gap, rank, auxiliary
restoration and step-down proofs remain valid and give \(p=R\).
Writing \(2n=R+\varepsilon+vE\), \(v\ge0\), a positive \(v\) would
force \(2n>2R\), contradicting \(n<p\). Thus \(2n=R+\varepsilon\).
If \(\varepsilon=-1\), then \(p=2n+1\); Pell duplication with
\(\chi_A(2)>P_0\) gives \(c>Ak>k(Y+1)\), contradicting the ratio
slacks. Therefore \(\varepsilon=\nu=1\), eliminating the exceptional
branch and proving all six equations. This proves the claimed equivalence
without using \(C>0\) prematurely.

\subsection{An explicit quartic in 33 positive unknowns}

The ordinary strong equation also removes one register from the
34-variable quadratic lift of package~08~\cite{src08}. Besides the
original fifteen positive variables, introduce eighteen positive
register variables
\[
 N,E,P,c,a,X,D,C_2,G,L,R,F_2,V_0,V,g,I,B_a,\Delta.
\]
In this subsection these letters are genuine unknowns. Only
\(q=\beta J+1\), \(k=\eta+\zeta\), \(H=4a+3\), and
\(C=q-F-\alpha\) are affine abbreviations. Impose the following
24 quadratic equations:
\begin{alignat}{2}\label{eq:quartic33-system}
 L&=q^2-1,&\qquad N&=q(L+1),\notag\\
 X&=wq,&a&=E+sN,\notag\\
 E&=X(a-E),&c&=k(a-E)+\eta,\notag\\
 P&=E(c-\eta)+k,&\tau^2&=P(P-k)+1,\notag\\
 D&=X+ac+\gamma H,&\Delta&=a^2+H,\notag\\
 C_2&=c^2,&D^2-\Delta C_2&=1,\notag\\
 k&=R+1+hE,&G&=(q-1)(q-F)+\alpha,\notag\\
 R&=LG+(M_C+qM_F)J,&\qquad(K+w)C+q-F-z(q-1)&=1,\notag\\
 F_2&=f^2,&V_0&=Tf-1,\notag\\
 V&=cV_0-RF_2,&g&=y^2-V^2,\notag\\
 I&=iC_2,&B_a&=I^2,\notag\\
 B_a&=\Delta(F_2-1),&y^2-B_ag&=1.
\end{alignat}
Every residual has degree at most two, since the four compiler ports are
fixed coefficients. If \(e_1,\ldots,e_{24}\) are these residuals, the
single polynomial is
\begin{equation}\label{eq:quartic33}
 \boxed{Q_{33}=\sum_{j=1}^{24}e_j^2=0.}
\end{equation}
Its degree is exactly four: the coefficient of \(\tau^4\) is 1.
The artifact supplies every expanded residual and the expanded quartic,
which has 224 formal terms before specializing its four coefficient ports.

The first equations recover \(N=q^3\), \(X=wq\), \(a-E=sq^3\),
\(E=X(a-E)\), and the original main Pell data. The last equations give
\(B_a=I^2=\Delta(f^2-1)\) and the ordinary auxiliary norm. Projection
therefore yields~\eqref{eq:ordinary-six}. Conversely, each ordinary
solution has the displayed register extension. All registers are
positive; the nontrivial cases are \(G>0\), \(V_0>0\), \(V>1\), and
\(g>0\), the last following from
\((W-1)(y^2-V^2)=V^2-1>0\). Thus positivity discards no solution.

The register saved relative to package~08 is its
\(K'=D^2-1=\Delta C_2\). That lift used \(I=iK'\), whereas the ordinary
version uses \(I=iC_2\). Deleting \(K'\) and replacing its two equations
by \(D^2-\Delta C_2=1\) gives precisely the one-variable, one-equation
reduction. Shifting all 33 coordinates by 1 gives the same degree-four
equation over nonnegative integers. These reductions inherit the stated
compiler and K215 calibration; they do not constitute a new formal
proof-assistant build or a proof of consistency.

